\documentclass[openany,fontset=fandol]{ctexbook}
\usepackage{amsmath,amssymb}
\usepackage[unicode,hidelinks]{hyperref}
\setlength{\emergencystretch}{2em}
\raggedbottom
\newcommand{\illis}{〔字迹不清〕}
\newcommand{\gap}{〔扫本缺损〕}
\begin{document}
\title{代数几何基础第五卷预备笔记\\手写提纲与异扫同文}
\author{据法文扫描及英译编辑稿校订}
\date{第87批\quad 2026年11月23日}
\maketitle

\chapter*{法文扫描第150页：手写提纲}

\noindent 以下按扫描的物理页码定位。上半页是未署日期的条目草记，下半页另有倒写算式与划改。
刊载机构的目录将物理第148--150页归在致迪厄多内的书信名下；
但第150页既没有承接前页落款的句子，也没有与英译稿第112页对应的正文，
因此作为邻接手稿单列，不添入书信。\footnote{参见格罗滕迪克资料库\href{https://webusers.imj-prg.fr/~leila.schneps/grothendieckcircle/EGA/EGAVcontents.pdf}{第五卷预备笔记目录}。该目录提供档案分组，本页手写项目不能直接视为书信提出的同一版目录，也不是正式刊本章节目录。}

\section*{正向条目}
\begin{description}
 \item[\S20] 关于 IV\illis{} 的补充；\illis{}；\illis{}，非双有理的\illis{}；以及爆破。\footnote{可辨首词为 \emph{complément}，后有形似罗马数字 IV 的标记，不能确定其后的编号。第二行可见 \emph{non} 及形如 \emph{birationnel} 的词根，末行可辨爆破一词；中间词和句法不强订。}
 \item[\S21] 除子，\illis{}，韦伊，皮卡尔；线性系，布劳尔--塞韦里。\footnote{在 \emph{Diviseurs} 后、\emph{Weil} 前有一组字迹可疑的缩记。此处不把它补成“卡蒂埃除子”，也不擅自指定“韦伊”“皮卡尔”两词的语法层级。}
 \item[\S22] 格拉斯曼簇（或概形）。\footnote{原文只写 \emph{Grassmanniennes}，未给出底概形或限定词。}
 \item[\S23] 超平面截面。
 \item[\S24] 无穷小扩张。\footnote{第23、24条分别为截面和无穷小扩张的缩记；第24条缩记不指定一个或多个扩张，本页无进一步定义。}
\end{description}

\section*{旋转半周后阅读的笔记}
可辨重划的一行、一个另行划去的算式、一个向右箭头，以及下方若干短线连成的草图。
可辨题词“对偶”，其限定语和下接词句\illis{}。
未划去的主行可作下列残读；缺字不是数学运算符：
\[
 H^{\text{〔疑为 }n\text{〕}}(X)\ \simeq\ \sum
       \text{〔指标及算符不清〕}(X,D^{\text{〔疑为 }i\text{〕}}).
\]
左、下另有简写与连线，字迹不足以判定两端表达式或箭头类型。\footnote{保留残读只是为了定位手稿墨迹。级数的指标、上标、项的首符号和划去文字都未能核定，不能从此行推出任何上同调分解或对偶定理。包内另附旋转后的局部原图，供回查，不作为重排正文。}

\section*{法文扫描第153页：卷首异扫}

这不是第四章第21节或书信的续页。它和法文扫描物理第1页
具有同一打字底文：下划线标题、16条拟定目录、目录后的编写说明及相同的半句截尾，
分行也基本对应。两页的边缘、反差和手写标记并不完全相同；不称为逐像素相同。
参考法文第1--2页及英译物理第2--3页可定位正文。
现行合订本已收相应目录与编写说明；以下供校勘，不计成新增一节。\footnote{第153页的原始编号栏多处不清。条号和部分词由法第1页及英第3页恢复；英第2页说明部分接续至法第2页。原档目录列出实存的第1--15节，草稿多列一条“改用线性系的语言”，二者不矛盾。}

\section*{超平面截面与锥形投影：拟定目录}
\begin{enumerate}
 \item 记号。
 \item 泛超平面截面：局部性质。
 \item 泛超平面截面：几何不可约性与几何连通性。
 \item 变动的超平面截面：“充分一般”的截面。
 \item 塞登伯格型定理。
 \item 任意超平面截面的连通性。
 \item 应用于构造特殊超平面截面和多截面。
 \item 例外超平面集合的维数。
 \item 改变射影浸入。
 \item 超平面截面铅笔与爆破后代数簇的纤维化。
 \item 格拉斯曼概形。
 \item 把前面的结果推广到线性截面。
 \item 初等态射与阿廷（M.~Artin）的一个定理。
 \item 锥形投影。
 \item 将前面的一些结果公理化。
 \item 改用线性系的语言表述。
\end{enumerate}

\section*{页末说明}
这一稿将最终稿里应当收入的各项结果混在一起，作了一份详细的提要。
要写成最终稿，就得把目前这个“第零阶段”的稿子彻底重组。
第一步大概是重新拟一份计划（现第11、12、14、15条想必会大幅提前）。
我还没有写第16条。原则上它不会有任何困难，也丝毫不用借助前面各条，
因为这里要做的\gap。\footnote{本页可读到“因为这里涉及……”一类残句，后半缺损。
法第1页同样断在此处，法第2页首行接为“不过是改换表述的问题”，英第2页也有对应句。
这里只在译注中交代接文，不据此声称第153页自身完整。法第2页其余文字已见合订本编写摘要。}

\section*{入册方式}
第150页作为书信相邻手稿保存，第153页作为卷首的复扫见证，从一份正文的方式收录。
法文卷末两个页位均已定位；这不等于法文171页逐页覆盖已验收，
也不等于全部上传卷册译完。下一步按物理页把法文、英译及现行译稿逐项对照；
例如须核定法第170--171页的边界，不机械沿用编辑目录的末页。
正式刊本 EGA IV.1 仍另从物理第221页提要及第223页（1.1.1）接续。

\end{document}
